\section*{Capítulo \ref{ch:g7:shadows-eclipses} --- Sombras y eclipses}

\begin{solution}{exo:g7:shadows-eclipses:1}
\omterm{prop:g7:shadows-eclipses:umbra}{Umbra}: la fuente está completamente escondida --- no llega
\omterm{def:g1:light-and-shadows:source}{luz} de ninguna parte de ella. \omterm{prop:g7:shadows-eclipses:umbra}{Penumbra}: la fuente está
escondida en parte --- sigue llegando \omterm{def:g1:light-and-shadows:source}{luz} de una parte de ella,
que asoma por encima del borde del obstáculo.
\end{solution}

\begin{solution}{exo:g7:shadows-eclipses:2}
Un punto no tiene ``partes'' que puedan quedar escondidas a medias:
cada sitio o lo ve o no lo ve --- \omterm{def:g1:light-and-shadows:shadow}{sombra} nítida, sin orla. Un cielo
cubierto es una fuente tan ancha que casi todos los sitios ven casi
toda ella: casi todo es \omterm{prop:g7:shadows-eclipses:umbra}{penumbra} y las \omterm{def:g1:light-and-shadows:shadow}{sombras} se
desdibujan.
\end{solution}

\begin{solution}{exo:g7:shadows-eclipses:3}
Los rayos que salen de los \emph{bordes} de la fuente y rozan los
lados del obstáculo. La \omterm{prop:g7:shadows-eclipses:umbra}{umbra} es la región a la que no llegan los
rayos de ninguno de los dos bordes.
\end{solution}

\begin{solution}{exo:g7:shadows-eclipses:4}
De Sol: Sol--Luna--Tierra, en ese orden --- la \omterm{def:g2:moon-in-the-sky:moon}{Luna} tiene que ser
nueva (entre nosotros y el Sol). De \omterm{def:g2:moon-in-the-sky:moon}{Luna}: Sol--Tierra--Luna --- la
\omterm{def:g2:moon-in-the-sky:moon}{Luna} tiene que estar llena (enfrente del Sol).
\end{solution}

\begin{solution}{exo:g7:shadows-eclipses:5}
La totalidad: solo quienes están dentro de la pequeña mancha que corre
de la \omterm{prop:g7:shadows-eclipses:umbra}{umbra}. Un mordisco parcial: los grandes alrededores de la
\omterm{prop:g7:shadows-eclipses:umbra}{penumbra}. Nada en absoluto: todos aquellos a quienes no alcanzan las
regiones de \omterm{def:g1:light-and-shadows:shadow}{sombra} --- casi todo el \omterm{def:g4:solar-system:planet}{planeta}.
\end{solution}

\begin{solution}{exo:g7:shadows-eclipses:6}
El \omterm{def:g7:shadows-eclipses:lunar}{eclipse de Luna} es la propia \omterm{def:g2:moon-in-the-sky:moon}{Luna} oscureciéndose --- un
cambio en un objeto que puede ver a la vez todo el lado nocturno. La
totalidad solar exige estar \emph{dentro} de la esbelta \omterm{prop:g7:shadows-eclipses:umbra}{umbra} de
la \omterm{def:g2:moon-in-the-sky:moon}{Luna} allí donde toca el suelo: una mancha, no un acontecimiento
del tamaño del cielo.
\end{solution}

\begin{solution}{exo:g7:shadows-eclipses:7}
La \omterm{def:g6:sun-earth-moon:axisorbit}{órbita} de la \omterm{def:g2:moon-in-the-sky:moon}{Luna} está inclinada unos cinco grados respecto
al plano Tierra-Sol: casi todas las lunas nuevas y
llenas pasan por encima o por debajo de la línea exacta, y las
\omterm{def:g1:light-and-shadows:shadow}{sombras} fallan. Solo cuando coinciden la fase y el cruce del plano
se alinean los tres mundos.
\end{solution}

\begin{solution}{exo:g7:shadows-eclipses:8}
La proyección con agujerito sobre un papel y los crecientes que un
árbol frondoso deja en el suelo (o un colador); y, para mirar
directamente, gafas de eclipse homologadas y sin daños. A ojo desnudo:
solo durante la totalidad misma, y la primera astilla que vuelve
termina el permiso.
\end{solution}

\begin{solution}{exo:g7:shadows-eclipses:9}
El Sol mide unas $400$ veces el ancho de la \omterm{def:g2:moon-in-the-sky:moon}{Luna} y está a unas
$400$ veces su distancia, así que los dos discos coinciden en el cielo
casi exactamente. Hijos de ese encaje: la totalidad existe, pero por
los pelos --- breve y con una manchita de \omterm{prop:g7:shadows-eclipses:umbra}{umbra} ---; y, cuando la
\omterm{def:g2:moon-in-the-sky:moon}{Luna} viaja por la parte lejana de su \omterm{def:g6:sun-earth-moon:axisorbit}{órbita} ovalada, su disco
algo más pequeño deja sin tapar el borde del Sol --- el anillo del
eclipse anular.
\end{solution}

\begin{solution}{exo:g7:shadows-eclipses:10}
La orla de \omterm{prop:g7:shadows-eclipses:umbra}{penumbra} de cualquier \omterm{def:g1:light-and-shadows:shadow}{sombra} la dibujan los bordes del
Sol, y la construcción enseña que la orla se ensancha con la distancia
entre el obstáculo y la pantalla: los rayos de los bordes opuestos del
Sol se separan más cuanto más viajan después del obstáculo. Los pies
están a distancia cero --- nítidos; la \omterm{def:g1:light-and-shadows:shadow}{sombra} de la cabeza queda a
dos \omterm{def:g2:measuring-length:units}{metros} de separación --- blanda.
\end{solution}

\begin{solution}{exo:g7:shadows-eclipses:11}
Cada hueco entre hojas es una cámara estenopeica natural; cada mancha
del suelo es la \omterm{prop:g5:mirrors-reflection:image}{imagen} del Sol que forma ese agujerito --- redonda
en los días corrientes. Durante el eclipse parcial, el propio Sol es
un creciente, así que cada camarita fiel proyecta un
creciente --- y todos apuntan hacia el mismo lado porque todos
copian al único creciente que hay en el cielo.
\end{solution}

\begin{solution}{exo:g7:shadows-eclipses:12}
La \omterm{def:g1:light-and-shadows:source}{luz} del Sol que roza el borde de la Tierra cruza la
\omterm{def:g5:air-pressure-weather:pressure}{atmósfera} por sus carreteras más largas posibles --- caminos de
longitud de \omterm{def:g1:day-and-night:daynight}{atardecer} en todo el contorno del \omterm{def:g4:solar-system:planet}{planeta} a la vez. El
\omterm{def:g2:air-around-us:air}{aire} difunde y se lleva buena parte de la \omterm{def:g1:light-and-shadows:source}{luz} que pasa (la
porción que pinta los cielos de la Tierra), y la \omterm{def:g7:light-sources-propagation:diffusion}{difusión} roba
antes los colores de camino corto, así que la \omterm{def:g1:light-and-shadows:source}{luz} que sobrevive al
largo viaje rasante y sigue hasta la \omterm{def:g2:moon-in-the-sky:moon}{Luna} es la superviviente del
\omterm{def:g1:day-and-night:daynight}{atardecer}: roja. La \omterm{def:g2:moon-in-the-sky:moon}{Luna} eclipsada brilla con la \omterm{def:g1:light-and-shadows:source}{luz} de
todos los \omterm{def:g1:day-and-night:daynight}{amaneceres} y \omterm{def:g1:day-and-night:daynight}{atardeceres} de la Tierra cayendo sobre
ella a la vez.
\end{solution}

\begin{solution}{pb:g7:shadows-eclipses:1}
\textbf{1.} La franja es el camino de la \omterm{prop:g7:shadows-eclipses:umbra}{umbra} de la \omterm{def:g2:moon-in-the-sky:moon}{Luna}
corriendo por el suelo; la ancha zona que la flanquea es la huella de
la \omterm{prop:g7:shadows-eclipses:umbra}{penumbra}.
\textbf{2.} Un eclipse parcial: el pueblo está en la
\omterm{prop:g7:shadows-eclipses:umbra}{penumbra} --- un mordisco en el Sol, mirado con la protección
adecuada, y nada de oscuridad.
\textbf{3.} Hasta una astilla de Sol del cinco por ciento es miles de
veces más brillante que la \omterm{ex:g2:moon-in-the-sky:shapes}{luna llena}: mantiene el cielo iluminado y
la corona invisible. La noche a mediodía pertenece únicamente a la
\omterm{prop:g7:shadows-eclipses:umbra}{umbra} --- no existe la ``casi totalidad''.
\textbf{4.} La mancha de \omterm{prop:g7:shadows-eclipses:umbra}{umbra} es más o menos redonda: un campamento
en la línea central queda bajo todo el ancho de la mancha mientras
esta pasa corriendo, mientras que los campamentos cercanos al borde de
la franja solo pillan una cuerda del círculo --- una travesía más
corta y una totalidad más corta.
\textbf{5.} Las gafas de eclipse homologadas bloquean la quemadura del
Sol; las gafas de sol no --- atenúan el deslumbramiento y dejan pasar
la parte que quema la retina, con lo que permiten quedarse mirando
cómodamente: peor que nada.
\textbf{6.} Agujerito en un extremo y pantalla de papel de calco en el
otro (o papel liso en el suelo): hay que colocarse \emph{de espaldas}
al Sol, con el agujero apuntando por encima del hombro, y mirar el
disco proyectado --- nunca a través del agujero.
\textbf{7.} Rechazada: el reflejo del agua es una
\omterm{prop:g5:mirrors-reflection:image}{imagen} especular del Sol, casi tan cegadora como el propio Sol ---
una quemadura rebotada sigue siendo una quemadura.
\textbf{8.} Decenas de imágenes en creciente sobre la sábana:
cada agujero es una cámara estenopeica y cada uno proyecta el Sol
mordido --- un colador lleno de eclipses.
\textbf{9.} La \omterm{def:g2:moon-in-the-sky:moon}{Luna} --- invisible contra el resplandor hasta que el
borde de su disco empieza a tapar el del Sol; el campamento ha entrado
en la \omterm{prop:g7:shadows-eclipses:umbra}{penumbra}.
\textbf{10.} El Sol visible se está encogiendo hasta quedar en una
astilla fina --- una fuente cada vez más puntual. Fuente más pequeña,
orlas de \omterm{prop:g7:shadows-eclipses:umbra}{penumbra} más estrechas: el borde blando de todas las
\omterm{def:g1:light-and-shadows:shadow}{sombras} se aprieta, y el mundo se vuelve extraño y nítido.
\textbf{11.} La propia \omterm{prop:g7:shadows-eclipses:umbra}{umbra} de la \omterm{def:g2:moon-in-the-sky:moon}{Luna} --- la clase está
dentro de la punta del cono de \omterm{def:g1:light-and-shadows:shadow}{sombra} mientras corre sobre el
campamento. Alrededor del disco negro cuelga la corona, la
\omterm{def:g5:air-pressure-weather:pressure}{atmósfera} exterior y nacarada del Sol, visible solo ahora; la
primera astilla del disco que vuelve restaura las reglas de la plena
\omterm{def:g1:light-and-shadows:source}{luz} del día --- gafas puestas.
\textbf{12.} Por ejemplo: ``Un eclipse es una \omterm{def:g1:light-and-shadows:shadow}{sombra} con mundos de
lámpara, de obstáculo y de pantalla: la \omterm{prop:g7:shadows-eclipses:umbra}{umbra} de la \omterm{def:g2:moon-in-the-sky:moon}{Luna}
rozando nuestro suelo, o la \omterm{def:g2:moon-in-the-sky:moon}{Luna} navegando por la nuestra ---
geometría que se puede dibujar con dos rayos y una regla, representada
a la mayor escala de la naturaleza''.
\end{solution}
